% No preamble. Prefix fmg-. Primary interface and proved finite geometry.
\section{Local gauge coordinates and the complete-density interface}
\label{fmg-section}

The homogeneous local trajectory does not yet identify all sectors
of the full Wilson integral. This section records a precise published
interface, proves a finite local gauge lemma, and shows why that
lemma alone does not supply a uniformly smooth physical background.

\subsection{The strongest published complete-density input}
\label{fmg-source}

Ba{\l}aban's Theorem~1 in \emph{Commun.\ Math.\ Phys.}\ 122
(1989), p.~355, supplies the complete effective-density class of
\emph{Commun.\ Math.\ Phys.}\ 119 (1988), Section~2, provided the
effective couplings of that construction remain sufficiently small.
The density bounds are (0.1), pp.~355--356; its large-field
exponentiation concludes at (1.98)--(1.102), pp.~390--391.
The original articles are available at
\href{https://scispace.com/pdf/large-field-renormalization-ii-localization-exponentiation-h8b7ih93xo.pdf}
{CMP122-II} and
\href{https://scispace.com/pdf/convergent-renormalization-expansions-for-lattice-gauge-29jyexd38e.pdf}
{CMP119}.

The relevant CMP119 input locations are:
\[
\begin{array}{c|l}
\text{pages/equations}&\text{interface}\\ \hline
243\text{--}244,\ (0.1)\text{--}(0.2)&
       \text{full pushforward and the additional operation}\\
246\text{--}248,\ (1.1)\text{--}(1.10)&
       \text{partitions and collars}\\
254\text{--}257,\ (2.1)\text{--}(2.18)&
       \text{multiscale determining sets and histories}\\
258,\ (2.19)\text{--}(2.22)&
       \text{large-field operator factorization}\\
258\text{--}262,\ (2.23)\text{--}(2.42)&
       \text{action, analytic domains and boundary terms}.
\end{array}
\]
In particular (2.38), p.~261, requires a local gauge with potential
and first-derivative bounds, alongside the other multiscale
regularity/current conditions. It is not a conclusion from one
bare plaquette threshold. The theorem concerns the specified
complete-density construction, not every small analytic local map.
OCR-damaged numerical constants in these source formulas are not
retranscribed here.

To use that result for our exact \(F_R=M_\ell^s\) local step,
one must either verify its complete induction class for this
operator or prove a replacement large-field theorem. In particular:
the full compact-field pushforward and partition must be specified;
the good-sector term must equal our actual retained-window
integral; exterior and transition data must be retained; and
normalized complementary-region operations must be bounded.
The source's own coupling recursion must be identified or replaced
by an independently justified one. Uniform weakness of our tuned
numbers alone is not equality of the two recursions.

Our strict product-ball families have an invariant total.
They do not automatically supply individually invariant local
orbit/current extensions or uniform analytic estimates on
multiscale determining sets. Nor does composition of averaging
constraints identify a one-shot window with sequential primitive
windows. These are concrete transfer obligations, without an
assertion that the source theorem cannot be adapted.

\subsection{An explicit finite cubical tree gauge}
\label{fmg-tree}

Let \(G\) be a compact matrix Lie subgroup of \(U(r)\), with
operator norm. Let
\(\Lambda_n=\{0,\ldots,n\}^d\), \(d\ge2\), carry its full cubical
cell structure. Oppositely oriented bonds carry inverse matrices.
For every vertex \(x\), let \(p(x)\) be the path from zero which
first increases coordinate one, then coordinate two, and so on.
These paths form a rooted spanning tree.

\begin{lemma}[Finite cubical gauge bound]
\label{fmg-tree-lemma}
Suppose every elementary plaquette obeys
\begin{equation}
                  \|U(\partial p)-I\|\le\epsilon.
 \label{fmg-plaquette-bound}
\end{equation}
Set
\begin{equation}
 v(x)=U(p(x)),\qquad
 U^v(x,y)=v(x)U(x,y)v(y)^{-1}.
 \label{fmg-gauge}
\end{equation}
Every tree edge of \(U^v\) is the identity, and each positive
\(\mu\)-edge satisfies
\begin{equation}
 \|U^v(x,x+e_\mu)-I\|
       \le\left(\sum_{\nu>\mu}x_\nu\right)\epsilon
       \le(d-1)n\epsilon.
 \label{fmg-edge-bound}
\end{equation}
Choose a fixed local exponential-chart radius
\(0<\delta_G\le1\). If \((d-1)n\epsilon<\delta_G\), write
\(U^v_b=\exp(i\eta A_b)\) in its unique local logarithm chart.
For physical mesh \(\eta\) and side \(\ell=n\eta\),
\begin{equation}
 \|A\|_\infty
       \le 2(d-1)n\epsilon/\eta
       =2(d-1)\ell\epsilon/\eta^2 .
 \label{fmg-physical-bound}
\end{equation}
\end{lemma}

\begin{proof}
The transformed edge is the product around the loop consisting
of \(p(x)\), the additional \(\mu\)-edge and reverse \(p(x+e_\mu)\).
Move the additional edge backwards past the later coordinate
segments. There are exactly \(\sum_{\nu>\mu}x_\nu\) elementary
square exchanges. Each exchange multiplies the loop holonomy
by a conjugate of a plaquette or its inverse, and the final
boundary word cancels. Unitary conjugation and inversion
preserve distance to \(I\). Moreover
\[
       \|W_1\cdots W_m-I\|
                     \le\sum_{j=1}^m\|W_j-I\|
\]
for unitary matrices, by telescoping. This proves
\eqref{fmg-edge-bound}; tree edges require no exchanges.
On the chosen local chart,
\(\|\log W\|\le2\|W-I\|\), as follows from the scalar eigenangle
bound. Dividing by \(\eta\) proves \eqref{fmg-physical-bound}.
\end{proof}

Before the logarithm, this gauge is a finite product of link
matrices and inverses. With a strict chart margin it is real
analytic and has a holomorphic continuation near each admitted
compact configuration. No refinement-uniform complex radius
or derivative bound follows just from
\eqref{fmg-edge-bound}. On a general simply connected cell complex
the same argument costs the filling areas of the tree fundamental
loops; simple connectivity alone does not bound those areas.

For a fixed finite current box the lemma gives local logarithmic
coordinates from sufficiently small plaquettes. For a fixed
physical box under refinement, its sufficient potential threshold
is \(\epsilon=O(\eta^2 b/\ell)\). It does not yet identify relative
fluctuations about a prescribed constrained minimum.

\subsection{A gauge-invariant obstruction to physical smoothness}
\label{fmg-holder-obstruction}

Fix \(0<\alpha<1\), a small \(f>0\), and a nonzero Cartan
generator \(H\) with \(\|H\|=1\). On a fixed physical cube
with decreasing mesh \(\eta\), let \(x\) denote integer lattice
coordinates and put
\begin{equation}
 U_1(x)=\exp\!\bigl(i\eta^2 f x_2\chi(x_1)H\bigr),\qquad
 U_\mu(x)=I\quad(\mu\ne1),
 \qquad
 \chi(k)=\begin{cases}0&k\text{ even},\\1&k\text{ odd}.
                     \end{cases}
 \label{fmg-oscillating-field}
\end{equation}
All non-\(12\) plaquettes are \(I\), and
\begin{equation}
               U(\partial p_{12}(x))
                    =\exp(-i\eta^2 f\chi(x_1)H).
 \label{fmg-oscillating-curvature}
\end{equation}
Thus every plaquette is within \(\eta^2f\) of \(I\).

\begin{proposition}[No uniform \(C^{1,\alpha}\) gauge]
\label{fmg-no-smooth-gauge}
For every fixed \(B<\infty\), sufficiently fine configurations
\eqref{fmg-oscillating-field} admit no gauge representative
\(U^v=\exp(i\eta A)\) on the cube satisfying
\begin{equation}
 \|A\|_\infty+\|\nabla_\eta A\|_\infty
                         +[\nabla_\eta A]_\alpha\le B.
 \label{fmg-smooth-assumption}
\end{equation}
The same assertion holds on an interior subcube containing
adjacent plaquettes of both parities.
\end{proposition}

\begin{proof}
For any \(A\) satisfying \eqref{fmg-smooth-assumption}, express
a \(\mu\nu\)-plaquette by the exact jet variables
\[
 a=A_\mu(x),\quad b=A_\nu(x),\quad
 c=\nabla_{\eta,\mu}A_\nu(x),\quad
 e=\nabla_{\eta,\nu}A_\mu(x).
\]
Its product is
\[
 P_\eta=
 e^{i\eta a}e^{i\eta b+i\eta^2c}
                  e^{-i\eta a-i\eta^2e}e^{-i\eta b}.
\]
The constant term is \(I\) and the first coefficient in \(\eta\)
vanishes. Consequently \((P_\eta-I)/\eta^2\) extends analytically
to \(\eta=0\), with uniformly bounded first jet derivatives
on each fixed bounded jet ball. Adjacent \(a,b\) differ by
at most \(B\eta\), and adjacent \(c,e\) differ by at most
\(B\eta^\alpha\). Hence
\begin{equation}
 \left\|(P_\eta(x)-I)/\eta^2
             -(P_\eta(x+e_\rho)-I)/\eta^2\right\|
                         \le C(B)\eta^\alpha .
 \label{fmg-curvature-holder}
\end{equation}
The same inequality holds for the difference of their norms.
But \(\eta^{-2}\|U(\partial p)-I\|\) is gauge invariant
at each plaquette. For \eqref{fmg-oscillating-curvature}
its adjacent values are zero and at least \(f/2\) for
sufficiently small \(\eta\). This contradicts
\eqref{fmg-curvature-holder}.
\end{proof}

This is an obstruction to every gauge, not merely to the chosen
tree gauge. The example is already in that tree gauge, where
\(A_1=\eta f x_2\chi(x_1)H\); its first derivative visibly grows
like \(f\ell/\eta\). The stronger invariant argument above
does not assert that a bounded \(C^1\) gauge alone is impossible.
It excludes the fixed physical \(C^{1,\alpha}\) background ball
needed here.

There is no contradiction with the complete-density method.
A raw fine field is not a uniformly smooth classical background.
The useful decomposition is a regular constrained background
plus rough fluctuations and explicitly retained exceptional
regions. A plaquette tail or a local logarithmic chart alone
does not construct that decomposition.

\subsection{Flat torus holonomies and transition functions}
\label{fmg-holonomy}

On a periodic direction of \(N\) bonds and physical circumference
\(\ell=N\eta\), consider
\[
       U_1(x)=e^{i\theta H/N},\qquad U_\mu(x)=I\quad(\mu\ne1).
\]
Its plaquettes are all \(I\), whereas the noncontractible loop is
\(e^{i\theta H}\), unchanged up to conjugacy by a periodic gauge.
If some periodic gauge has \(U_b^v=e^{i\eta A_b}\) with
\(\|A\|_\infty\le b\), telescoping the loop gives
\begin{equation}
                  \|e^{i\theta H}-I\|\le N\eta b=\ell b.
 \label{fmg-holonomy-bound}
\end{equation}
Thus zero plaquette curvature cannot force an arbitrarily small
global physical potential. Local tree gauges trivialize this
flat field on contractible cubes, but their transition data
retain the holonomy. On overlaps the exact relation is
\begin{equation}
 U^{(i)}(x,y)=h_{ij}(x)U^{(j)}(x,y)h_{ij}(y)^{-1},
 \qquad h_{ij}=v_i v_j^{-1}.
 \label{fmg-transitions}
\end{equation}
Those data must be retained or integrated consistently. They
cannot be set to identity because local plaquettes are small.

\subsection{A scale-comparison fact supplied by local shooting}
\label{fmg-couplings}

Our actual weak finite trajectories obey
\(\gamma_j=g_j^{-2}\ge\Gamma_{\mathrm w}>0\) and
\(|\gamma_{j+1}-\gamma_j|\le D_{\mathrm{abs}}\).
Telescoping yields, for any two stages,
\begin{equation}
 \frac{\gamma_i}{\gamma_j}
       \le1+\frac{D_{\mathrm{abs}}}{\Gamma_{\mathrm w}}|i-j|,
 \qquad
 \frac{g_i}{g_j}
       \le\left(1+\frac{D_{\mathrm{abs}}}
                          {\Gamma_{\mathrm w}}|i-j|\right)^{1/2}.
 \label{fmg-coupling-comparison}
\end{equation}
For the second inequality apply the first with \(i,j\) interchanged.
Both reverse ratios satisfy the same bounds. No one-step
monotonicity is used. This is a useful polynomial interscale
comparison; it is not asserted to match every source radius
inequality without checking its constants and domains.

\subsection{The next normalized interface}
\label{fmg-next}

A direct full-measure route needs an exact gauge-invariant
good/bad partition and a normalized conditional representation.
Its good component must be the actual stationary-plus-fluctuation
integral with its actual retained window. Its complementary
replacement must be summable in component size, scale age and
boundary geometry, uniformly in the allowed exterior data.
A fixed-box conditional Wilson estimate is a genuine first input;
its constants need not survive growing components or a nonlinear
average constraint. Conversely, a normalized estimate on a
coarse-average fiber cannot be inferred by conditioning an
unconditional tail without controlling that fiber's normalization.

The import route makes the same obligations through the published
complete-density class instead of a new direct proof. Neither
route has been completed by the local trajectory or by the
geometric lemmas in this section. No continuum-existence conclusion
is asserted here.
